\documentclass[10pt,a4paper]{amsart}
\usepackage[margin=19mm]{geometry}
\usepackage{amsmath,amssymb,mathtools}
\usepackage[scr=rsfs,cal=euler]{mathalfa}
\usepackage{xfrac}
\usepackage[unicode,hidelinks,bookmarks=false]{hyperref}
\newcommand{\ie}{i.e.\ }
\newcommand{\cf}{cf.\ }
\newcommand{\resp}{respectively\ }
\newcommand{\Subsection}{\subsection}
\providecommand{\nobreakdash}{\nobreak\hbox{-}}
\setlength{\emergencystretch}{2em}
\multlinegap0pt
\usepackage{bm}
\usepackage{bbm}
\usepackage{dsfont}
\usepackage{xspace}
\usepackage{cancel}
\usepackage{mathtools}
\usepackage{amsthm}
\makeatletter
\@ifundefined{claim}{\newtheorem{claim}{Claim}}{}
\@ifundefined{theorem}{\newtheorem{theorem}{Theorem}}{}
\@ifundefined{lemma}{\newtheorem{lemma}[theorem]{Lemma}}{}
\makeatother
\makeatletter
\DeclareMathOperator{\Reg}{Reg}
\expandafter\ifx\csname itemizenew\endcsname\relax
\newenvironment{itemizenew}{\begin{itemize}[leftmargin=2pc]}{\end{itemize}}
\fi
\makeatother
\title[Compactness in abelian group theory]{Compactness in abelian group theory}
\author{Filippo Calderoni, Ava Ostrem}
\date{Source: arXiv:2508.16038v2 (CC BY 4.0)}
\begin{document}
\maketitle

\noindent\emph{Source and licence.} Compactness in abelian group theory. Version arXiv:2508.16038v2 (CC BY 4.0), distributed under the Creative Commons Attribution 4.0 International licence, as recorded in the arXiv licence metadata for this exact version and shown on the abstract page https://arxiv.org/abs/2508.16038v2. Full rights notice: licenses/arxiv-2508.16038-CC-BY-4.0.txt.\par\smallskip

\noindent\textit{Editorial scope.} Counted unit: the Theorem whose source label is \texttt{thm:compactness sigma cyclic} in the article (source0.tex lines 536--539, source label \texttt{thm:compactness sigma cyclic}), delivered together with the article's own complete original proof of it (source0.tex lines 541--580, one \texttt{proof} environment of 40 lines). No mathematics is written, repaired or re-derived: statement and proof are the article's own text line for line. Required context retained uncounted: thm : weak comp Sigma cyclic (lines 200--203); lem:reflection (lines 382--386); thm:Gamma invariant groups (lines 496--499). Every definition, display and label of those lines is retained unchanged. Omitted from this package and not redistributed: 1--199, 204--381, 387--495, 500--535, 540--540, 581--672. Nothing inside a retained excerpt is omitted or altered: every retained body is delivered exactly as the article's lines are written, so no omission receipt of any kind accompanies this package. Citations: the retained excerpts carry no citation command at all, so no bibliography entry of the article is transcribed and no reference is redistributed. Reference adaptation: none was needed and none was made; every reference inside the delivered unit resolves inside it, so no unresolved reference or citation is produced. Printed slips kept literal: no printed word, symbol, hypothesis or number of the counted statement or of its proof was altered. Layout and class: the article's own class is \texttt{amsart} with packages including fontenc, mathpazo, eucal, microtype, hyperref, amssymb, enumitem, mathrsfs, mathtools, tikz, upgreek, verbatim, hyphenat, url; the wrapper is the lane's pinned \texttt{amsart} editorial wrapper at 10pt on A4, which restates the theorem-like environments the retained text needs and adds the article's own macro definitions that the retained text needs (\texttt{\textbackslash Reg}), with extra wrapper packages (bm, bbm, dsfont, xspace, cancel, mathtools). The delivered numbering is the unit's own. Assets: no retained span contains a figure environment, a graphics-inclusion command, a PDF-page inclusion, a table or a TikZ picture; no image file is copied into this package, the package declares no asset dependency and no image file is redistributed with it. Licence: version arXiv:2508.16038v2 of this article is distributed under the Creative Commons Attribution 4.0 International licence, as recorded in the arXiv licence metadata for this exact version and shown on the abstract page https://arxiv.org/abs/2508.16038v2. A full rights notice is delivered at licenses/arxiv-2508.16038-CC-BY-4.0.txt. Characters: every retained excerpt is pure ASCII, so the delivered engine, XeLaTeX with its default Latin Modern text font, reports no missing character.
\begin{theorem}
\label{thm : weak comp Sigma cyclic}
     If \(\kappa\) is weakly compact, then every \(\kappa\)-\(\Sigma\)-free abelian group of cardinality \(\kappa\) is \(\Sigma\)-cyclic.
\end{theorem}

\begin{lemma}
    \label{lem:reflection}
        Let \(\kappa\) be a weakly compact cardinal and let \(\{S_\alpha: \alpha<\kappa\}\) be a family of stationary subsets of \(\kappa\). Then there is a stationary \(T\subseteq \kappa \cap \Reg(\kappa)\) such that for all \(\lambda \in T\) and \(\alpha<\lambda\), the set
        \( S_\alpha \cap \lambda\) is stationary in \(\lambda\).
    \end{lemma}

    \begin{theorem}
        \label{thm:Gamma invariant groups}
        Let \(A\) be \(\kappa\)-\(\Sigma\)-cyclic with \(|A|=\kappa\). Then \(A\) is \(\Sigma\)-cyclic if and only\footnote{We denote by \(0\) the equivalence class \([\emptyset]_\sim\).} if \(\Gamma(A) = 0\).
    \end{theorem}

\begin{theorem}
        \label{thm:compactness sigma cyclic}
        Let \(\kappa\) be weakly compact. If \(A\) is an abelian group with \(|A| = \kappa\) and \(A\) is \(\kappa\)-\(\Sigma\)-cyclic, then \(A\) is \(\Sigma\)-cyclic.
    \end{theorem}

    \begin{proof}[Proof of Theorem~\ref{thm : weak comp Sigma cyclic}]
    Let \(\kappa\) be weakly compact. Assume that \(A\) is an abelian group with \(|A| = \kappa\) and that \(A\) is \(\kappa\)-\(\Sigma\)-cyclic.

    
    Suppose \(A\) is not \(\Sigma\)-cyclic towards contradiction.
        Let \(\{A_\alpha: \alpha<\kappa\}\) be a filtration for \(A\) consisting of \(\Sigma\)-cyclic groups.
        Without loss of generality, assume that \(A_0 = \{0\}\) and \(|A_\alpha|\leq{|\alpha|+\aleph_0}\). Then let
        \[
        E = \{\alpha<\kappa :
        \{ \beta>\alpha : \beta<\kappa \text{ and } A_\beta/A_\alpha\text{ is not \(\Sigma\)-cyclic}\}
        \text{ is stationary in \(\kappa\)}.
        \}
        \]
        Since \(A\) is not \(\Sigma\)-cyclic, \(E\) is stationary by Theorem~\ref{thm:Gamma invariant groups}. Now we build a family \(\{S_\alpha\}\) of stationary subsets of \(\kappa\)  by setting
        \[
        S_\alpha =
        \begin{cases}
            \{\beta>\alpha : \beta<\kappa \text{ and } A_\beta/A_\alpha \text{ is not \(\Sigma\)-cyclic}\} &\alpha \in E.\\
            E &\text{otherwise}.
        \end{cases}
        \]
        Let \(T\subseteq \kappa \cap \Reg(\kappa)\) be a stationary set such that for all \(\lambda\in T\) and \(\alpha<\lambda\) the set \(S_\alpha\cap \lambda\) is stationary in \(\lambda\).

        Now pick any \(\lambda \in T\). First note that \(E\cap \lambda\) is stationary in \(\lambda\). This is because 
        \[\{\beta>0 : A_\beta / A_0\text{ is not \(\Sigma\)-cyclic}\} = \emptyset\] is clearly not stationary. Therefore, \(0 \notin E\) and consequently
        \(S_0\cap \lambda = E\cap \lambda\), which is stationary by the assumption on \(T\). Also, note that \(\{A_\alpha: \alpha<\lambda\}\) is a filtration for the \(\Sigma\)-cyclic group \(A_\lambda\).

        \begin{claim}For any \(\lambda \in T\), we have
            \[E\cap \lambda \subseteq \{ \alpha<\lambda : \{ \beta>\alpha: \beta<\lambda \text{ and \(A_\beta/A_\alpha\) is not \(\Sigma\)-cyclic}\} \text{ is stationary in } \lambda \}\] and both sets are stationary in $\lambda$.
        \end{claim}

        \begin{proof}
            For \(\alpha \in E\cap \lambda\), we know that \(S_\alpha\) reflects at \(\lambda\) by Lemma~\ref{lem:reflection} and the choice of \(T\). This means that the set
            \[S_\alpha \cap \lambda  = \{\beta > \alpha : \beta <\lambda \text{ and } A_\beta/A_\alpha \text{ is not \(\Sigma\)-cyclic}\}
            \]
            is stationary in \(\lambda\). Since \(E\cap \lambda\) is stationary, so is the larger set.
        \end{proof}

         Given the claim, \(\Gamma(A_\lambda)\neq 0\). This shows that \(A_\lambda\) is not \(\Sigma\)-cyclic due to Theorem~\ref{thm:Gamma invariant groups}, and this is a contradiction.
    \end{proof}

\end{document}
